%% Question bank chapter file - generated by bank_tools.py build (do not edit while the bank is still being built:
%% it is regenerated from the per-test data in Question Bank\data). Every question is one bankq block:
%%   \begin{bankq}[id=..., source={...}, paper=.., qno=.., type=..., subtopic={...}, answer={...}, answer_src=key|solved,
%%                 working={...}, difficulty=1-10, diff_parts={K S I C F}, diff_note={...},
%%                 group=<passage id or empty>, also_in={duplicate copies in other tests}]
%% difficulty = min(10, K+S+I+C+F-1): K concepts 1-3, S steps 1-4, I insight 0-2, C computation 0-2, F format 0-1
%%   ...question LaTeX, options, figures (\bfig / \bfigi -> img/)...
%%   \end{bankq}
%% Passages shared by several questions are bankpassage blocks; their questions carry group=<passage id>.
\documentclass[10pt]{article}
\usepackage[a4paper,margin=1.8cm]{geometry}
\usepackage[T1]{fontenc}
\usepackage{amsmath,amssymb,mathtools}
\usepackage{graphicx}
\usepackage[export]{adjustbox}
\usepackage{xcolor,needspace,keyval}
\usepackage{enumitem}
\usepackage{booktabs,array,tabularx}
\usepackage[hidelinks]{hyperref}
\setlength{\parindent}{0pt}
\setlength{\parskip}{3pt}
\makeatletter
\define@key{bq}{id}{\def\bq@id{#1}}\define@key{bq}{source}{\def\bq@src{#1}}\define@key{bq}{paper}{\def\bq@paper{#1}}
\define@key{bq}{qno}{\def\bq@qno{#1}}\define@key{bq}{type}{\def\bq@type{#1}}\define@key{bq}{subtopic}{\def\bq@sub{#1}}
\define@key{bq}{answer}{\def\bq@ans{#1}}\define@key{bq}{answer_src}{\def\bq@asrc{#1}}\define@key{bq}{working}{\def\bq@work{#1}}
\define@key{bq}{group}{\def\bq@group{#1}}\define@key{bq}{also_in}{\def\bq@also{#1}}
\define@key{bq}{difficulty}{\def\bq@diff{#1}}\define@key{bq}{diff_parts}{\def\bq@dparts{#1}}\define@key{bq}{diff_note}{\def\bq@dnote{#1}}
\newenvironment{bankq}[1][]{\def\bq@work{}\def\bq@also{}\def\bq@group{}\def\bq@diff{?}\def\bq@dparts{}\def\bq@dnote{}\setkeys{bq}{#1}%
  \par\needspace{4\baselineskip}\medskip\noindent{\small\textbf{\bq@id}\hfill\color{gray}[\bq@src, Paper \bq@paper, Q\bq@qno]
  \bq@type\ $\cdot$ \bq@sub\ $\cdot$ \textbf{D\bq@diff}}\par\nopagebreak\relax}
  {\par{\small\color{teal}Answer: \textbf{\bq@ans}\ifx\bq@work\@empty\else\ (solved: \bq@work)\fi
  \ifx\bq@also\@empty\else\par\color{gray}Also in: \bq@also\fi}\par}
\define@key{bp}{id}{\def\bp@id{#1}}\define@key{bp}{source}{\def\bp@src{#1}}\define@key{bp}{questions}{\def\bp@qs{#1}}
\newenvironment{bankpassage}[1][]{\setkeys{bp}{#1}\par\needspace{4\baselineskip}\bigskip\noindent
  {\small\textbf{Passage \bp@id}\hfill\color{gray}[\bp@src; questions \bp@qs]}\par\nopagebreak\relax\itshape}{\par}
\makeatother
\newcommand{\bfig}[2]{\par\smallskip{\centering\includegraphics[width=#1pt,max width=0.62\linewidth,max height=0.32\textheight]{img/#2}\par}\smallskip}
\newcommand{\bfigi}[2]{\includegraphics[width=#1pt,max width=\linewidth,max height=0.32\textheight,valign=c]{img/#2}}
\newcommand{\bmissing}[1]{\textit{[figure missing in the source conversion: \detokenize{#1}]}}
\newenvironment{qoptions}{\begin{itemize}[leftmargin=1.8em,itemsep=1pt,topsep=2pt,label={}]}{\end{itemize}}
\newcommand{\opt}[2]{\item[(#1)] #2}
\providecommand{\fwstart}{}\providecommand{\fwend}{}\providecommand{\qfullsplit}{}
\newcolumntype{L}[1]{>{\raggedright\arraybackslash}p{#1}}
\newcolumntype{C}[1]{>{\centering\arraybackslash}p{#1}}
\newcolumntype{R}[1]{>{\raggedleft\arraybackslash}p{#1}}
\providecommand{\sectionhdr}[1]{}\providecommand{\subsectionhdr}[1]{}


\begin{document}
{\Large\bfseries s-block --- Question Bank}\par\smallskip
{\small Chemistry chapter C35. Questions from all converted test papers, grouped by source series. Answers: printed key where it matches the paper, otherwise solved (answer\_src=solved, with working). Exact repeats across papers appear once, with also\_in listing the other copies.}\par\medskip\hrule

\section*{TOAS papers}
\begin{bankq}[id=TOAS01-P2-CHEM-S1-Q11, source={TOAS - Test 01 (30 Jan 2022)}, paper=2, qno=11, type=single, subtopic={Alkali metals}, answer={C}, answer_src=key, difficulty=1, diff_parts={K1 S1 I0 C0 F0}, diff_note={concepts: trends of inorganic compounds | steps: colour intensity and decomposition temperature; solubility and thermal stability orders}]
Which of the following is not a correct order\\
(i) AgF $<$ AgCl $<$ AgBr (Colour intensity)\\
(ii) $\mathrm{Rb_2CO_3>K_2CO_3>Na_2CO_3}$ (decomposition temperature)\\
(iii) $\mathrm{BeC_2O_4<MgC_2O_4<CaC_2O_4}$ (Solubility)\\
(iv) $\mathrm{Li_3N>Na_3N>K_3N}$ (Thermal stability)
\begin{qoptions}
\opt{A}{i, iii, iv}
\opt{B}{i, iii}
\opt{C}{only iii}
\opt{D}{iii and iv}
\end{qoptions}
\end{bankq}
\begin{bankq}[id=TOAS02-P1-CHEM-S1-Q05, source={TOAS - Test 02 (13 Feb 2022)}, paper=1, qno=5, type=single, subtopic={Alkali metals}, answer={C}, answer_src=key, difficulty=1, diff_parts={K1 S1 I0 C0 F0}, diff_note={concepts: periodic trends of alkali metals | steps: trend in each plotted property; match against the species}]
Which of the following graph correctly depicts the property mentioned against the indicated species.
\begin{qoptions}
\opt{A}{\bfigi{128.6}{TOAS02_p1_chem_s1_q05_A.png}}
\opt{B}{\bfigi{132.5}{TOAS02_p1_chem_s1_q05_B.png}}
\opt{C}{\bfigi{131.8}{TOAS02_p1_chem_s1_q05_C.png}}
\opt{D}{\bfigi{144.1}{TOAS02_p1_chem_s1_q05_D.png}}
\end{qoptions}
\end{bankq}
\begin{bankq}[id=TOAS05-P2-CHEM-S1-Q19, source={TOAS - Test 05 (03 Apr 2022)}, paper=2, qno=19, type=single, subtopic={Alkali metals}, answer={C}, answer_src=key, difficulty=1, diff_parts={K1 S1 I0 C0 F0}, diff_note={concepts: trends of inorganic compounds | steps: colour and decomposition temperature; solubility and thermal stability}]
Which of the following is not a correct order\\
(i) $\mathrm{AlCl_3<AlBr_3<AlI_3}$ (Covalent character)\\
(ii) $\mathrm{Rb_2CO_3>K_2CO_3>Na_2CO_3}$ (decomposition temperature)\\
(iii) $\mathrm{BeC_2O_4<MgC_2O_4<CaC_2O_4}$ (Solubility)\\
(iv) $\mathrm{Li_3N>Na_3N>K_3N}$ (Thermal stability)
\begin{qoptions}
\opt{A}{i, iii, iv}
\opt{B}{i, iii}
\opt{C}{only iii}
\opt{D}{iii and iv}
\end{qoptions}
\end{bankq}
\begin{bankq}[id=TOAS08-P2-CHEM-S3-Q02, source={TOAS - Test 08 (15 May 2022)}, paper=2, qno=2, type=numerical, subtopic={Alkaline earth metals}, answer={6.00}, answer_src=key, difficulty=2, diff_parts={K1 S1 I0 C0 F1}, diff_note={concepts: hydrate of magnesium chloride | steps: MgCl2.6H2O crystallises}]
During extraction of Magnesium by Dow's process, magnesium hydroxide is treated with HCl followed by concentration of solution and crystallisation to give crystals of $\mathrm{MgCl_2.xH_2O}$. The value of ``$x$'' is \underline{\hspace{1.5cm}}.
\end{bankq}
\begin{bankq}[id=TOAS15-P1-CHEM-S1-Q06, source={TOAS - Test 15 (22 Aug 2022)}, paper=1, qno=6, type=single, subtopic={Alkali metals}, answer={A}, answer_src=key, difficulty=1, diff_parts={K1 S1 I0 C0 F0}, diff_note={concepts: density trend of alkali metals | steps: potassium is the only alkali metal lighter than sodium}]
Which of the following having density less than Na?
\begin{qoptions}
\opt{A}{K}
\opt{B}{Rb}
\opt{C}{Cs}
\opt{D}{None of these}
\end{qoptions}
\end{bankq}
\begin{bankq}[id=TOAS17-P2-CHEM-S1-Q19, source={TOAS - Test 17 (18 Sep 2022)}, paper=2, qno=19, type=single, subtopic={Alkali metals}, answer={A}, answer_src=key, difficulty=2, diff_parts={K1 S2 I0 C0 F0}, diff_note={concepts: alkali metal compounds | steps: solubility in ethanol; oxides with oxygen; fluoride and oxide solubility}]
Which of the following statements are correct ?
\par\noindent(a) Both LiCl and MgCl$_2$ are soluble in ethanol.\par
\noindent(b) The oxides Li$_2$O and MgO combine with excess of oxygen to give superoxide.\par
\noindent(c) LiF is less soluble in water than other alkali metal fluorides.\par
\noindent(d) Li$_2$O is more soluble in water than other alkali metal oxides.\par
Choose the most appropriate answer from the options given below:
\begin{qoptions}
\opt{A}{(a) and (c) only}
\opt{B}{(a), (c) and (d) only}
\opt{C}{(b) and (c) only}
\opt{D}{(a) and (c) only}
\end{qoptions}
\end{bankq}
\begin{bankq}[id=TOAS17-P2-CHEM-S2-Q07, source={TOAS - Test 17 (18 Sep 2022)}, paper=2, qno=7, type=single, subtopic={Alkali metals}, answer={321}, answer_src=solved, working={P: Na, Li and Cs give the normal oxide with limited air, so 3; Q: Li, Ca and Sr chlorides colour the flame red, so 2; R: Li, Be and Al chlorides are covalent, so 1}, difficulty=3, diff_parts={K1 S2 I1 C0 F0}, diff_note={concepts: properties of alkali metals and the diagonal elements | steps: oxide formed in limited air; flame colours; covalent chlorides}]
Match given properties with elements
\begin{center}\footnotesize\setlength{\tabcolsep}{2pt}
\begin{tabular}{|L{0.35cm}L{2.9cm}|L{0.35cm}L{1.4cm}|}\hline
\multicolumn{2}{|c|}{List-I (Properties)}&\multicolumn{2}{c|}{List-II (Elements)}\\\hline
P&Produce stable M$_2$O type oxide with limited supply of air&1&Li, Be, Al\\\hline
Q&Chloride of elements produce red colour in oxidising flame&2&Li, Ca, Sr\\\hline
R&Chlorides of such elements having significant covalent character or covalent nature in anhydrous form&3&Na, Li, Cs\\\hline
\end{tabular}\end{center}
Match list-I and list-II and give your answer as abc where a,b and c are the match of 1, 2, 3 from list-II.

\textbf{(for example Matching of P is = 2, Q = 3 \& R = 1 then your answer is 231)}
\end{bankq}

\section*{TMAS papers}
\begin{bankq}[id=TMAS04-P2-CHEM-S1-Q06, source={TMAS - Test 04 (19 Nov 2023)}, paper=2, qno=6, type=multi, subtopic={Alkali metals}, answer={A,C,D}, answer_src=key, difficulty=3, diff_parts={K1 S1 I0 C1 F1}, diff_note={concepts: alkali metal properties | steps: check each}]
Alkali metals are characterised by their :
\begin{qoptions}
\opt{A}{high electropositive character}
\opt{B}{high reduction potentials}
\opt{C}{low melting points}
\opt{D}{high solubility in liquid ammonia at $-33^\circ$C.}
\end{qoptions}
\end{bankq}
\begin{bankq}[id=TMAS04-P2-CHEM-S3-Q05, source={TMAS - Test 04 (19 Nov 2023)}, paper=2, qno=5, type=numerical, subtopic={Alkaline earth metals}, answer={4}, answer_src=key, difficulty=3, diff_parts={K1 S1 I0 C1 F1}, diff_note={concepts: hardness of water | steps: classify each salt; count}]
How many aqueous solution of compounds cause hardness in water?

$\mathrm{NaCl}$, $\mathrm{CaCl_2}$, $\mathrm{CaSO_4}$, $\mathrm{Mg(HCO_3)_2}$, $\mathrm{KNO_3}$, $\mathrm{Mg(NO_3)_2}$
\end{bankq}

\section*{TNAS papers}
\begin{bankq}[id=TNAS04-P2-CHEM-S1-Q05, source={TNAS - Test 04 (18 Jun 2023)}, paper=2, qno=5, type=multi, subtopic={Alkaline earth metals}, answer={A,D}, answer_src=key, difficulty=5, diff_parts={K1 S2 I1 C1 F1}, diff_note={concepts: properties of alkali and alkaline earth compounds | steps: thermal stability trend; solubility trend; melting point trend; Fajans rule}]
Choose the \textbf{CORRECT} order for the given properties
\begin{qoptions}
\opt{A}{$\mathrm{MgSO_4<SrSO_4<BaSO_4}$ : Thermal stability order}
\opt{B}{$\mathrm{BeC_2O_4<CaC_2O_4<BaC_2O_4}$ : Solubility in water}
\opt{C}{$\mathrm{LiCl>NaCl>KCl}$ : Melting point order}
\opt{D}{$\mathrm{BeF_2>CaF_2>SrF_2}$ : Covalent character order}
\end{qoptions}
\end{bankq}

\section*{Star Review (Other) papers}
\begin{bankq}[id=STAR03-P1-CHEM-S2-Q09, source={Star Review (Other) - Test 03 (04 Dec 2023)}, paper=1, qno=9, type=multi, subtopic={Alkali metals}, answer={A,B,C,D}, answer_src=key, difficulty=4, diff_parts={K1 S2 I1 C0 F1}, diff_note={concepts: properties of alkali metal compounds | steps: reaction of ammonium chloride; Solvay process; solutions in ammonia; fire extinguisher}]
Which of the following statement(s) is/are correct?
\begin{qoptions}
\opt{A}{When NaOH is added to $\mathrm{NH_4Cl}$ then ammonia is evolved}
\opt{B}{$\mathrm{K_2CO_3}$ cannot be prepared by Solvay process}
\opt{C}{The alkali metals dissolve in liquid $\mathrm{NH_3}$ giving deep blue solution which are conducting in nature}
\opt{D}{$\mathrm{NaHCO_3}$ with $\mathrm{H_2SO_4}$ is used in fire extinguisher}
\end{qoptions}
\end{bankq}

\section*{Nurture (General) papers}
\begin{bankq}[id=NUR01-P1-CHEM-S3-Q01, source={Nurture (General) - Test 01 (12 Jun 2021)}, paper=1, qno=1, type=single, subtopic={Alkaline earth metals}, answer={5}, answer_src=solved, working={B is CO2 (C = +4) and D is C2H2 (C = -1); the difference is 5}, difficulty=5, diff_parts={K2 S2 I2 C0 F0}, diff_note={concepts: identification of gases from limestone and calcium carbide, alkaline earth metals | steps: identify B and D; oxidation states; difference}]
Lime stone(X) $\xrightarrow{\Delta}$ Solid (A) + gas (B)\par
Other Binary Compound of Calcium(Y) $\xrightarrow{\mathrm{H_2O}}$ Solution (C) + Gas (D)\par
Gas (B) turns solution (C) milky but does not react with $\mathrm{KMnO_4/H^+}$.\par
Gas (D) produce white ppt. in tollen reagent and it also produce benzene when passed through red hot iron tube.\par
The difference of oxidation state of same central atom in (B) \& (D) is :
\end{bankq}
\begin{bankq}[id=NUR01-P1-CHEM-S3-Q03, source={Nurture (General) - Test 01 (12 Jun 2021)}, paper=1, qno=3, type=single, subtopic={Alkali metals}, answer={0}, answer_src=solved, working={x=3 (K, Rb, Cs), y=2 (Be, Mg), u=5, so 3+2-5=0}, difficulty=4, diff_parts={K2 S2 I1 C0 F0}, diff_note={concepts: alkali and alkaline earth metal facts, alkali metals | steps: superoxide formers; hydrolysis of hydrated halides; ammonia solutions; combine}]
Number of non-radioactive alkali metals forming superoxide as major product on heating with excess $\mathrm{O_2}=x$,\par
Number of II-A metals of periodic table whose hydrated halides suffer hydrolysis on heating $=y$.\par
Number of non radio active alkali metals which dissolve in liquid ammonia and produce blue colour solution $=u$\par
The value of $(x+y-u)$ would be....
\end{bankq}
\end{document}
